% Intro and Related Work section draft for CHO manuscript
% Consumed by writing-agent in Stage 3.

\section{Introduction}

Neural oscillations in the mammalian brain are thought to coordinate activity across distributed regions and to support the integration of sensory input, the control of movement, and the maintenance of cognitive state \citep{buzsaki2004neuronal,fries2015rhythms,giraud2012cortical,jensen2010shaping}. Detecting these oscillations is therefore a foundational step in much of systems and cognitive neuroscience: their presence, duration, location, and fundamental frequency each carry information about the underlying neural process. In practical terms, three questions are routinely asked of a recording. \emph{When} does an oscillation occur, i.e. what are its onset and offset times? \emph{Where} is it generated, i.e. which brain region or layer gives rise to it? And \emph{what} is its fundamental frequency, i.e. which band does it truly occupy \citep{buzsaki2004neuronal,pfurtscheller1999eventrelated}? Answering the \emph{when}, \emph{where}, and \emph{what} is a prerequisite for more ambitious questions about the \emph{why}, the \emph{how}, and the \emph{whom} (the pathology) of neural rhythms.

The dominant analytical tool for answering these questions is the Fast Fourier Transform (FFT) and its derivatives. FFT-based methods decompose a signal into sinusoidal components and then search for peaks over a $1/f$ aperiodic background, a strategy made rigorous by the specparam/FOOOF framework \citep{donoghue2020parameterizing} and extended to time-frequency by SPRiNT \citep{wilson2022sprint} and to discrete oscillation events by OEvent \citep{neymotin2022oevent}. These methods assume, implicitly or explicitly, that the underlying oscillations are sinusoidal. A growing body of work, however, shows that neural oscillations are routinely non-sinusoidal: mu rhythms, hippocampal theta, and cortical beta all exhibit systematic asymmetries between peaks and troughs \citep{belluscio2012crossfrequency,cole2017waveform,cole2019cycle,donoghue2022methodological}. When an FFT is applied to a non-sinusoidal rhythm, the asymmetry of the waveform is expressed as a cascade of harmonic peaks at integer multiples of the fundamental frequency. These harmonic peaks are artifacts of the sinusoidal basis, yet they appear indistinguishable from independent oscillations in the power spectrum. The result is an inflated false-positive detection rate that contaminates downstream analyses and, in some cases, overturns physiological interpretation.

The ambiguity is concrete. An unfiltered electrocorticographic recording from auditory cortex can display peaks over $1/f$ at 6, 12, and 18 Hz. Without further analysis, it is impossible to tell whether these are three independent rhythms or a single 6 Hz non-sinusoidal oscillation with two harmonics. Phase-phase coupling analysis \citep{belluscio2012crossfrequency} can resolve the question: if the 18 Hz component is phase-locked to the 6 Hz component in a 1:3 ratio, it is a harmonic. But the permutation tests required by this approach scale poorly across channel counts, trial counts, and frequency pairs. Phase-phase coupling is therefore impractical for the real-time and closed-loop settings in which oscillation detection is most needed, such as brain-computer interfaces, neurofeedback systems, and phase-locked deep brain stimulation \citep{cagnan2017phase,schalk2004bci2000}.

In this study, we define the principle criteria that characterize a neural oscillation and synthesize them into an algorithm that simultaneously recovers the \emph{when}, the \emph{where}, and the \emph{what} of non-sinusoidal oscillations. The method, which we call Cyclic Homogeneous Oscillation (CHO) detection, combines three ideas. First, following FOOOF, CHO uses a parametric $1/f$ estimate to threshold the time-frequency map and generate candidate bounding boxes---refining the empirical thresholding of OEvent \citep{neymotin2022oevent}. Second, CHO rejects any candidate that exhibits fewer than two full cycles, thereby excluding evoked and event-related potentials whose spectral signatures mimic theta or alpha oscillations \citep{decheveigne2019filters}. Third, CHO tests each surviving bounding box with auto-correlation: the positive-peak spacing of the auto-correlation of the raw signal within the box must agree with the center frequency of the box. Because auto-correlation reads periodicity directly from the time-domain signal, it is insensitive to waveform shape, and the test rejects the harmonic peaks that plague FFT-based approaches. A box whose center frequency is 24 Hz but whose raw signal periodicity is 7 Hz is rejected. We evaluate CHO on synthetic non-sinusoidal bursts (2.5 cycles, 1--3 s long) convolved with $1/f$ noise, and on four human datasets: ECoG from 8 subjects and scalp EEG from 7 subjects during a pre-stimulus auditory reaction-time task, resting-state ECoG (6 subjects), and resting-state SEEG from 6 subjects with depth electrodes in the hippocampus. Across these datasets we compare CHO to FOOOF \citep{donoghue2020parameterizing}, OEvent \citep{neymotin2022oevent}, and SPRiNT \citep{wilson2022sprint}.

\section{Related Work}
\label{sec:relwork}

\paragraph{Frequency-domain detection.} The FFT and its descendants separate a recording into sinusoidal components within canonical bands. Peaks over a $1/f$ aperiodic floor are interpreted as oscillations. FOOOF (now specparam) formalized this by jointly fitting a Lorentzian-like aperiodic component and a mixture of Gaussian peaks \citep{donoghue2020parameterizing}, and it is the de facto standard for peak-frequency estimation. Its principal limitation in our setting is that it operates only on the average spectrum and has no native mechanism for rejecting harmonic peaks.

\paragraph{Time-frequency and event-based detection.} SPRiNT extends FOOOF into the time-frequency domain by parameterizing short-time spectra \citep{wilson2022sprint}. OEvent identifies bounding boxes of elevated power in a wavelet map as discrete oscillation events and summarizes each by center frequency, duration, and number of cycles \citep{neymotin2022oevent}. Neither approach, however, tests whether the bounding box corresponds to a genuine fundamental frequency or to a harmonic of a lower-band rhythm.

\paragraph{Time-domain and waveform-based detection.} A separate strand of work studies oscillations directly in the time domain, quantifying peak-to-trough asymmetry, rise-decay asymmetry, and instantaneous frequency \citep{cole2017waveform,cole2019cycle,quinn2021within,gips2017lineclock}. These methods are well-suited to describing waveform shape \emph{after} an oscillation has been identified, but they do not themselves adjudicate whether a candidate spectral peak is fundamental or harmonic.

\paragraph{Cross-frequency coupling.} Belluscio et al. \citep{belluscio2012crossfrequency} established that $n{:}m$ phase-phase coupling can identify harmonic relations, and follow-up work has used this tool to disentangle theta-gamma and theta-beta couplings in hippocampus and cortex. The method is statistically principled but computationally heavy, requiring surrogate distributions across all channel-frequency pairs, which puts it out of reach of online analysis pipelines.

\paragraph{The gap.} No existing method jointly (i) detects the onset and offset of a rhythm, (ii) rejects harmonic peaks without resorting to permutation-based cross-frequency tests, and (iii) scales to real-time analysis suitable for closed-loop neuromodulation. CHO is designed to fill this gap by combining bounding-box detection with an auto-correlation periodicity test, at a computational cost that is dominated by a single FFT and auto-correlation per candidate box.
